% !TeX spellcheck = es_ES
\documentclass[letterpaper,11pt]{article}
\usepackage[utf8]{inputenc}
\usepackage[T1]{fontenc}
\usepackage[spanish]{babel}
\decimalpoint % separador decimal con punto, como en las tablas de la fuente original
\usepackage{lmodern,microtype}
\usepackage{amsmath,amssymb,amsfonts}
\usepackage{geometry}
\geometry{margin=2.54cm}
\usepackage{setspace}
\setstretch{1.15}
\usepackage{enumitem}
\usepackage{booktabs}
\usepackage{color}
\usepackage{hyperref}
\definecolor{ntr}{rgb}{0.8,0,0}
\hypersetup{colorlinks=true,linkcolor=ntr,citecolor=ntr,urlcolor=black}

% Portada (Luis Cabral): titulo \LARGE sans-serif alineado a la izquierda, sin numero de pagina
\makeatletter
    \renewcommand{\@maketitle}{%
    \newpage\parindent0in\null\vskip2em%
    {\LARGE\textbf{\textsf{\@title}}\par\vskip2em}%
    \@author\par\vskip2em%
    \@date%
    \par
    \thispagestyle{empty}\setcounter{page}{1}%
    }%
\makeatother

\DeclareMathOperator{\Var}{Var}
\DeclareMathOperator{\Cov}{Cov}
\DeclareMathOperator{\Corr}{Corr}
\DeclareMathOperator{\E}{\mathbb{E}}
\DeclareMathOperator{\N}{Normal}
\DeclareMathOperator{\LP}{\mathbb{L}}
\newcommand\independent{\protect\mathpalette{\protect\independenT}{\perp}}
\def\independenT#1#2{\mathrel{\rlap{$#1#2$}\mkern2mu{#1#2}}}
% Teoria asintotica
\newcommand{\SR}{\mathbb{R}}
\DeclareMathOperator{\plim}{plim}
\DeclareMathOperator{\Avar}{AVar}
\newcommand{\pto}{\stackrel{p}{\longrightarrow}}
\newcommand{\dto}{\stackrel{d}{\longrightarrow}}
\newcommand{\adistr}{\stackrel{a}{\sim}}
\newcommand{\epsi}{\varepsilon}


\begin{document}

	\title {Análisis Econométrico  \\ Tarea 5}
	\author{Magíster en Economía, Universidad Alberto Hurtado}
	\date{Fecha de entrega: 26 de noviembre de 2026 (al principio de la ayudantía)}

\maketitle

Total: 100 puntos

Los ejercicios 2 y 4 requieren estimar modelos con datos. Utilice R (u otro software
estadístico) y adjunte su código. Reporte errores estándar robustos a
heteroscedasticidad y utilice un nivel de significancia del 5\% salvo que se indique otra
cosa. Los datos están en el repositorio del curso, en formato R:
\href{https://github.com/ramirode/econometria_mae/raw/main/problem-sets/datos/card.rds}{\texttt{card.rds}} y
\href{https://github.com/ramirode/econometria_mae/raw/main/problem-sets/datos/jec.rds}{\texttt{jec.rds}}.
Se leen con \texttt{readRDS}.

\begin{enumerate}

\item (20 puntos) Considere el modelo con una variable explicativa endógena $y_2$
    \begin{align*}
    y_{1}=z_{1}'\delta_1+\alpha\, y_{2}+u,
    \end{align*}
    con $\E(zu)=0$, donde $z=(z_1',z_2')'$, $z_1$ incluye la constante y las demás variables
    explicativas exógenas, y $z_2$ son los instrumentos. Considere el siguiente procedimiento
    para estimar el modelo:
    \begin{enumerate}
        \item Regresar $y_2$ en una constante y $z_2$ (\emph{no} en $z_1$) y obtener los valores
            predichos $\hat y_2=\hat\pi_0+\hat\pi'z_2$.
        \item Regresar $y_1$ en $z_1$ y $\hat y_2$ para obtener $\hat\delta_1$ y $\hat\alpha$.
    \end{enumerate}
    Demuestre que, en general, $\hat\delta_1$ y $\hat\alpha$ son inconsistentes. ¿Cuándo serán
    consistentes?

    \emph{Pista:} sea $y_2^0=\pi_0+\pi'z_2$ la regresión poblacional de $y_2$ en una constante
    y $z_2$, y sea $a_2$ el error de proyección, de modo que $y_2=y_2^0+a_2$ con $\E(a_2)=0$ y
    $\E(a_2z_2)=0$. Para simplificar, suponga que $\pi_0$ y $\pi$ son conocidos, es decir,
    asuma que $\hat y_2$ es $y_2^0$. Escriba
    \begin{align*}
    y_1=z_1'\delta_1+\alpha\, y_2^0+u+\alpha\, a_2,
    \end{align*}
    y verifique si el error compuesto está correlacionado con las variables explicativas.

\item (30 puntos) Utilice los datos de \texttt{card} del trabajo de Card (1993) sobre los
    retornos de la educación: 3,010 hombres de la encuesta \textit{National Longitudinal
    Survey of Young Men}, con datos de 1976. Las variables son $lwage$ (log del salario por
    hora), $educ$ (años de educación), $exper$ (años de experiencia), $black$, $south$ y $smsa$
    (dummies para raza negra, vivir en el sur y vivir en un área metropolitana en 1976),
    $smsa66$ (dummy para vivir en un área metropolitana en 1966), $reg661,\dots,reg669$
    (dummies de región de residencia en 1966), $nearc4$ y $nearc2$ (dummies iguales a 1 si en
    1966 vivía cerca de una universidad de cuatro y de dos años, respectivamente) e $iq$
    (puntaje de un test de IQ, disponible para una parte de la muestra). Considere el modelo
    \begin{align*}
    \log(wage_i)=\beta_0&+\beta_1\, educ_i+\beta_2\, exper_i+\beta_3\, exper_i^2+\beta_4\, black_i
    +\beta_5\, south_i\\&+\beta_6\, smsa_i+\beta_7\, smsa66_i+\boldsymbol\beta_{reg}'\, reg_i+u_i,
    \end{align*}
    donde $reg=(reg661,\dots,reg668)'$.
    \begin{enumerate}
        \item (6 puntos) Estime el modelo por MCO. Construya un intervalo de confianza al 95\%
            para el retorno de la educación $\beta_1$.
        \item (6 puntos) Card propone como instrumento de $educ$ la variable $nearc4$. Estime la
            primera etapa, incluyendo los controles del modelo. Contraste el supuesto de
            relevancia.
        \item (6 puntos) Estime el modelo por variables instrumentales con $nearc4$ como
            instrumento. Compare el intervalo de confianza al 95\% para $\beta_1$ con el de (a).
            ¿Qué explica la diferencia en el ancho de los intervalos?
        \item (6 puntos) Estime el modelo por 2SLS con $nearc4$ y $nearc2$ como instrumentos.
            Contraste la relevancia conjunta de los dos instrumentos y realice el contraste de
            sobreidentificación de Sargan. ¿Qué concluye? Compare con el resultado de (c).
        \item (6 puntos) Regrese $iq$ en $nearc4$. ¿Está $iq$ correlacionado con $nearc4$?
            Regrese ahora $iq$ en $nearc4$, $smsa66$, $reg661$, $reg662$ y $reg669$. ¿Están $iq$
            y $nearc4$ parcialmente correlacionados? ¿Qué concluye sobre la importancia de
            controlar por la ubicación geográfica en 1966 para la validez de $nearc4$ como
            instrumento?
    \end{enumerate}

\item (20 puntos) La demanda de un bien está dada por $q_i^d=\alpha_0+\alpha_1 p_i+u_i$, donde
    $q_i^d$ denota la cantidad demandada, $p_i$ el precio y $u_i$ otros factores distintos del
    precio que afectan la cantidad demandada. La oferta del bien está dada por
    $q_i^s=\beta_0+\beta_1 p_i+v_i$, donde $q_i^s$ denota la cantidad ofrecida y $v_i$ otros
    factores distintos del precio que afectan la cantidad ofrecida. Suponga que $\E(u_i)=0$,
    $\E(v_i)=0$, $\Cov(u_i,v_i)=0$ y $\alpha_1\neq\beta_1$. Las cantidades observadas son de
    equilibrio, por lo tanto $q_i=q_i^s=q_i^d$.
    \begin{enumerate}
        \item (5 puntos) Resuelva el sistema de ecuaciones para mostrar cómo $q_i$ y $p_i$
            dependen de $u_i$ y $v_i$.
        \item (5 puntos) Calcule $\Cov(p_i,u_i)$ y $\Cov(p_i,v_i)$.
        \item (5 puntos) Tenemos una muestra aleatoria de $q_i$ y $p_i$ y regresamos $q_i$ en
            $p_i$ (con constante). ¿A qué converge en probabilidad el estimador de la pendiente,
            $\hat\alpha_1$, cuando se lo interpreta como estimador de la pendiente de la
            demanda? ¿Cuál es el signo del sesgo asintótico?\\
            \emph{Pista:} utilice $\Cov(p_i,u_i)$ y la ecuación de demanda.
        \item (5 puntos) Interprete ahora la misma pendiente como estimador de la pendiente de
            la oferta, $\hat\beta_1$. ¿A qué converge en probabilidad? ¿Cuál es el signo del
            sesgo asintótico?\\
            \emph{Pista:} utilice $\Cov(p_i,v_i)$ y la ecuación de oferta.
    \end{enumerate}
    En (c) y (d) suponga que la demanda tiene pendiente negativa y la oferta pendiente positiva.

\item (30 puntos) Utilice los datos de \texttt{jec} del trabajo de Porter (1983, \textit{Bell
    Journal of Economics}) para estimar la función de demanda de transporte ferroviario de
    granos. A partir de 1880 un cartel conocido como \textit{Joint Executive Committee} (JEC)
    controló el transporte ferroviario de granos desde el Medio Oeste hacia las ciudades del
    este de EE.UU. El cartel era anterior a la ley antimonopolio de 1890 y por lo tanto operaba
    legalmente. Cada tanto, los miembros del cartel se desviaban del precio acordado, lo que
    provocaba un colapso temporal del acuerdo colusivo. En este ejercicio utilizamos las
    variaciones en la oferta asociadas al funcionamiento del cartel para estimar la elasticidad
    de la demanda de transporte ferroviario de granos. Los datos son semanales (328 semanas).
    Suponga que la demanda es
    \begin{align*}
    \log(Q_t)=\beta_0+\beta_1\log(P_t)+\beta_2\, ice_t+\sum_{j=2}^{12}\beta_{1+j}\, seas_{j,t}+u_t,
    \end{align*}
    donde $Q_t$ son las toneladas de granos transportadas en la semana $t$, $P_t$ es el precio
    de transportar una tonelada de granos por ferrocarril, $ice_t$ es una dummy igual a 1 si los
    Grandes Lagos no son navegables por el hielo, y $seas_{j,t}$ es una dummy igual a 1 si la
    semana $t$ pertenece al mes $j$ (el mes 1 es la categoría omitida). La variable $cartel_t$
    es una dummy igual a 1 si en la semana $t$ el cartel estaba operando (es decir, no había
    colapsado).
    \begin{enumerate}
        \item (6 puntos) Estime la ecuación de demanda por MCO. Interprete el resultado.
        \item (6 puntos) Explique por qué la interacción entre oferta y demanda puede hacer
            inconsistente la estimación por MCO. Sugiera el signo del sesgo asintótico.
        \item (6 puntos) Considere utilizar $cartel_t$ como instrumento de $\log(P_t)$. Utilice el
            razonamiento económico para justificar que es un instrumento válido.
        \item (6 puntos) Estime la primera etapa. Contraste el supuesto de relevancia.
        \item (6 puntos) Estime la ecuación de demanda por variables instrumentales. Interprete
            los resultados y compare con los de MCO. Utilice el intervalo de confianza al 95\%
            para evaluar si la elasticidad estimada es consistente con el comportamiento de un
            cartel que maximiza beneficios.
    \end{enumerate}

\end{enumerate}
\end{document}
